\Article{11}
To illustrate this method further, we give one example from among its many applications. The experiments carried out on $1832$ Sept. $11$ yielded the following figures.
\SourceLeaf{102}
\expandafter\def\csname HistoricalNoteText10\endcsname{For example, the amplitude of the first needle's oscillations in the first experiment was $0^0\ 37'26''$ at the beginning and $0^0\ 28'34''$ at the end. In the fifth experiment it was $1^0\ 10'21''$ at the beginning. After $177$ oscillations it was $0^0\ 45'35''$ and after $677$ oscillations it was $0^0\ 6'44''$.}
\begin{center}
\begin{tabular}{c||c|c||c}
&\multicolumn{3}{c}{Simultaneous oscillations}\\
Exp.&\multicolumn{2}{c||}{first needle}&second needle\\\cline{1-4}
&Load&one oscillation&one oscillation\\
I&$r=180^{\mathrm{mm}}$&$24''63956$&$17''32191$\\
II&$r=130$&$20{,}77576$&$17{,}32051$\\
III&$r=80$&$17{,}66798$&$17{,}31653$\\
IV&$r=30$&$15{,}80310$&$17{,}30529$\\
V&unloaded&$15{,}22990$&$17{,}31107$
\end{tabular}
\end{center}

The times were measured by a chronometer that lost $14''24$ in a mean-time day. Each of the two weights $p$ was $103{,}2572$ grams. The distances $r$ were determined in millimeters with microscopic precision. The duration of one oscillation was inferred from at least $100$ oscillations (in the fifth experiment, even from $677$ for the first needle) and has already been reduced to the duration of an infinitesimal oscillation. In any case, these corrections are imperceptible because of the very small oscillation amplitude\HistoricalNote{10} that our apparatus allows us to use while retaining the highest precision. We first reduce these oscillation times to the mean value of $TM$ that prevailed during the fifth experiment, following the instructions in the preceding article under II. We then reduce them to the values that would have been obtained without torsion, by multiplying by $\BareRoot{\frac{n+1}{n}}$, where $n$ is $=424{,}8$ in the first four experiments and $=597{,}4$ in the fifth (see Article $9$), and finally to mean solar time by multiplying by $\frac{86400}{86385{,}76}$. In this way we obtain
\begin{center}
\begin{tabular}{crl}
I.&$24''65717=t'$&for $r'=180^{\mathrm{mm}}$\\
II.&$20{,}79228=t''$&for $r''=130^{\mathrm{mm}}$\\
III.&$17{,}68610=t'''$&for $r'''=80^{\mathrm{mm}}$\\
IV.&$15{,}82958=t''''$&for $r''''=30^{\mathrm{mm}}$\\
V.&$15.24515=t$&for the unloaded needle.
\end{tabular}
\end{center}

Taking the second, the millimeter and the milligram as the units of time, distance and mass, so that $p=103257.2$, we obtain the following from the first and fourth experiments.
\SourceLeaf{103}
\[
TM=179641070,\qquad K+C=4374976000
\]
and then, from the fifth experiment,
\[
K=4230282000,\quad \text{as well as}\quad C=144694000
\]
If, however, all the experiments are to be included in the calculation, the method of least squares is most conveniently applied as follows. We start from the approximate values of the unknowns $x,y$ obtained by combining the first and fourth experiments. Denoting the corrections still to be added by $\xi,\eta$, we set
\[
\begin{aligned}
x&=88{,}13646+\xi\\
y&=21184{,}85+\eta
\end{aligned}
\]
In this way the calculated values of the times $t'.\ t'',t''',t''''$ are obtained by straightforward methods
\[
\begin{aligned}
t'&=24{,}65717-0{,}13988\xi+0{,}00023008\eta\\
t''&=20{,}78731-0{,}11793\xi+0{,}00027291\eta\\
t'''&=17{,}69121-0{,}10036\xi+0{,}00032067\eta\\
t''''&=15{,}82958-0{,}08980\xi+0{,}00035838\eta
\end{aligned}
\]
Comparing them with the observed values by the method of least squares yields
\[
\begin{aligned}
\xi&=-0{,}03230,&\quad\eta&=-12{,}38\\
x&=88{,}10416,&y&=21172{,}47
\end{aligned}
\]

From this we finally obtain
\[
TM=179575250,\qquad K+C=4372419000
\]
and then, using the first experiment,
\[
K=4228732400,\qquad C=143686600
\]

Here is a comparison of the times calculated from the corrected values of $x,y$ with the observed times.
\begin{center}
\begin{tabular}{c|r|r|r}
Exp.&Calculated time&Observed time&Difference\\\hline
I&$24''65884$&$24''65717$&$+0''00167$\\
II&$20{,}78774$&$20{,}79228$&$-0{,}00454$\\
III&$17{,}69046$&$17{,}68610$&$+0{,}00436$\\
IV&$15{,}82805$&$15{,}82958$&$-0{,}00153$
\end{tabular}
\end{center}

\SourceLeaf{104}
We take the length of the simple pendulum at Göttingen to be $=994^{\mathrm{mm}}126$. Gravity, measured in the unit of accelerating forces underlying the preceding calculations, is therefore $=9811{,}63$. If we prefer to take gravity itself as the unit, we obtain $TM=18302{,}29$. This number expresses how many milligrams must press under gravity on a lever one millimeter long to produce the same effect as the force with which terrestrial magnetism tends to turn that needle about the vertical axis.

\Article{12}
Having finished determining the product of the horizontal terrestrial magnetic force $T$ and the magnetic moment $M$ of a given needle, we now turn to the second part of the investigation, the determination of the quotient $\frac MT$. We will obtain it by comparing the action of this needle on another needle with that of terrestrial magnetism on the same needle. As explained in the introduction, the comparison can be made either in motion or in equilibrium. We have often tried both methods. Since the latter is far preferable to the former for several reasons, we will confine our investigation here to it. This is all the more appropriate since the former can be treated without difficulty in exactly the same way.

\Article{13}
The principle of virtual displacements makes it very easy to reduce the conditions for equilibrium of a movable body acted on by arbitrary forces to a single formula. For each force, take the infinitesimal displacement of its point of application, project it in the direction of that force and multiply it by the force. The sum of these products must have no positive value for any virtual displacement, that is, any displacement compatible with the general conditions to which the body's motion is subject. Hence, insofar as virtual displacements in opposite directions are possible everywhere, that sum, which we will denote by $\mathrm{d}\Omega$, must be $=0$ for every virtual displacement.

The movable body we consider here is a magnetic needle whose point $G$ is attached to a twistable thread fixed at its upper end. The thread merely prevents the distance from $G$ to the fixed end of the thread from exceeding the length of the thread. Thus here too, as in the case of a completely free body, the position of the body in space depends on six variables, and its equilibrium accordingly on six conditions. But since the solution of the problem here need only serve to determine the quotient $\frac MT$, \SourceLeaf{105}it suffices to consider the virtual displacement that consists in rotation about the vertical axis through $G$. Clearly, we may regard this axis as fixed and take as the sole variable the angle between the vertical plane containing the needle's magnetic axis and the plane of the magnetic meridian. We will measure this angle from the north part of the meridian toward the east and denote it by $u$.

